\subsection{配比--损失关系与跨体系域映射}
配比实验给出 $512$ 组 $17$ 维领域配比与对应的 $13$ 域验证损失。借鉴按配比回归损失的思路\cite{regmix2024}，以配比向量 $\bm p$ 为特征、分域损失为目标，对每个有损失记录的目标域 $k$ 建立对数域线性模型；配比输入满足单纯形约束：
\begin{equation}
\ln \hat L_k(\bm p)=a_k+\sum_{i=1}^{17} b_{ki}\,p_i,
\qquad \text{s.t.}\quad p_i\ge 0,\ \sum_{i=1}^{17}p_i=1 .
\label{eq:q1-mix}
\end{equation}
由于 $\sum_i p_i=1$，截距和全部配比系数不能同时唯一识别。计算时把系数改写为 $\sum_i b_{ki}=0$ 的等价零和编码，并相应平移截距，使预测不变。图中的平均绝对系数只表示相对平均域的线性对比强度，不是因果影响。配比输入本身归一化到单纯形；另用梯度提升回归作拟合对照。

\textbf{求解与验证。}线性模型在 $13$ 个目标域上的训练决定系数介于 $0.50$--$0.79$（中位约 $0.67$）。表~\ref{tab:q1-mixture-holdout} 把三个留出尺度的两种指标并列：同尺度 $1$m 的平均 $R^2$ 为 $0.664$；在 $60$m 与 $1$B 上，平均 $R^2$ 分别降至 $-5.387$ 和 $-265.290$，说明固定的 $1$m 截距与系数不能直接预测更大模型的绝对损失。平均 Spearman 相关仍为正，表明域内配比排序保留一定信息，但相关系数不能保证选出的最优配比在新尺度上仍最优。梯度提升模型在 $1$m 数据的五折交叉验证平均 $R^2$ 为 $0.977$；它只检验同尺度拟合，不验证跨尺度泛化。

\begin{table}[H]
  \centering
  \caption{配比回归的分尺度留出检验}
  \label{tab:q1-mixture-holdout}
  \small
  \begin{tabular}{lrrr}
    \toprule
    检验尺度 & 配比数 & 平均 $\ln L$ 的 $R^2$ & 平均 Spearman 相关 \\
    \midrule
    $1$m  & 256 & $0.664$ & $0.844$ \\
    $60$m & 256 & $-5.387$ & $0.842$ \\
    $1$B  & 64  & $-265.290$ & $0.736$ \\
    \bottomrule
  \end{tabular}
\end{table}

上述数值均对有验证损失的 $13$ 个目标域求平均。$R^2<0$ 表示预测平方误差大于直接使用该检验集目标域平均 $\ln L$ 的误差，不能仅以“尺度不同”略过。若把该模型用于新规模，应先用该规模的少量已知配比重新标定截距，再用独立配比验证排序和最优选择的遗憾值；本文没有进行这种再标定。现有排序相关只是一条迁移线索，零和编码的系数对比见图~\ref{fig:q1_mixture_imp}。

\begin{figure}[H]
  \centering
  \includegraphics[width=0.9\textwidth,height=0.8\textheight,keepaspectratio]{fig_q1_mixture_importance.pdf}
  \caption{零和编码下配比系数的相对对比强度}
  \label{fig:q1_mixture_imp}
\end{figure}

图~\ref{fig:q1_mixture_imp} 按零和编码后的平均绝对系数排序。即使采用可复现的编码，域间相关仍会改变系数；这张图只展示线性模型内的相对对比，当前问题二的广义律也未使用这组系数。

跨体系域映射方面，质量信号覆盖的域与配比涉及的 $17$ 个域并不一一对应。按提供的映射指南，其中三个域可直接同名对应、三个域语义近似对应，其余多数为无同名质量域的推断域。对推断域按体裁相似性借用最接近的质量信号（如数学类借用 arxiv 与数学信号），这一按重要性加权借用的思路与基于重要性重采样的数据选择一致\cite{dsir2023}，并标注其可信度较低；四个既无对应质量域又无损失列的域（如邮件、议会记录等）仍保留其配比作为式~\eqref{eq:q1-mix} 的解释变量，间接影响其它域损失，而非无声剔除。这样既显式处理了体系差异，又不丢失任何配比信息。

\subsection{小结}
本章得到域级综合质量（中位数 $0.079$ 用作问题三的 $Q_0$）和约 $4.96\%$ 的局部双高冲突。配比回归在跨尺度检验中保留了一定排序相关，但绝对损失预测显著失准；其系数可用于后续配比建模，尚未进入本文的广义律或预算优化数值计算。
